(1) 压力满足 \(\partial_t^2p=c^2\nabla^2p\)。代入轴对称谐波得 \(P''+r^{-1}P'+\kappa^2P=0\)，其中 \(\kappa^2=\omega^2/c^2-k^2\)。圆心正则解为 \(P=A J_0(\kappa r)\)，舍去在 \(r=0\) 对数发散的 \(Y_0\)。径向特征量代表横截面驻波结构，不能直接当作轴向传播波数。

(2) 线性声学动量方程 \(\rho_0\partial_t v_r=-\partial_rp\)，刚壁不透流 \(v_r(a)=0\) 等价于 \(\partial_rP(a)=0\)。因此 \(J_1(\kappa a)=0\)，取非负根 \(\xi_n\)（含 \(\xi_0=0\)），有 \(\kappa_n=\xi_n/a\)，轴向色散关系 \(\omega^2=c^2[k^2+\xi_n^2/a^2]\)。若采用固定压强的软壁则边界改为 \(J_0(\kappa a)=0\)，两者不可混淆。

(3) 基模 \(\xi_0=0\) 没有横向梯度，频率 \(\omega=c|k|\)，代表管内平面声波，无正截止频率。第一个非零 \(J_1\) 根为 \(\xi_1\approx3.83170597\)，对应第一轴对称横向激发的截止角频率 \(\omega_{\rm cut}=c\xi_1/a\)。若 \(\omega<\omega_{\rm cut}\)，该模的 \(k\) 变成纯虚数，沿管道指数衰减而非携能传播。

(4) 对高于截止的模 \(k>0\)，色散给 \(v_p=\omega/k=c\sqrt{1+\xi_n^2/(a^2k^2)}>c\)，群速度 \(v_g=d\omega/dk=c^2k/\omega<c\)，故严格乘积 \(v_pv_g=c^2\)。相速度超声速并不表示信号传递超过声速，波包能量速度由群速度控制。临近截止 \(k\to0\) 时相速度发散而群速度趋零，体现波导模式的截止。
